\section{Implementierung}

\subsection{Methodisches Vorgehen}

Die Implementierung folgte einem \textbf{issue-basierten Ansatz}:
Jede Funktion wurde als separates GitHub-Issue spezifiziert,
in einem Feature-Branch implementiert, getestet und per Merge
in \texttt{develop} integriert.

Die Issues waren in drei Phasen gegliedert:
\begin{itemize}
  \item \textbf{Phase 1}: Proto-Definition, Projektstruktur
  \item \textbf{Phase 2}: Dispatcher-Implementierung (Issues \#13--\#21)
  \item \textbf{Phase 3}: Containerisierung mit Docker (Issues \#34--\#38)
\end{itemize}

\subsection{Dispatcher-Implementierung (Issues \#13--\#21)}

Die Dispatcher-Implementierung wurde vollständig abgeschlossen.
Alle 49 Integrationstests bestehen.
Die wichtigsten implementierten Funktionen sind in
Abschnitt~\ref{sec:technologien} detailliert beschrieben.

Übersicht der implementierten Issues:

\begin{table}[H]
\centering
\begin{tabular}{lll}
\toprule
\textbf{Issue} & \textbf{Titel} & \textbf{Hauptdatei} \\
\midrule
\#13 & POST\_TASK (Aufgabe einreichen)       & \texttt{servicer.py} \\
\#14 & LOOKUP\_WORKER (Namensdienst-Abfrage) & \texttt{namensdienst\_client.py} \\
\#15 & DISPATCH (Task an Worker senden)      & \texttt{dispatch\_loop.py} \\
\#16 & RESULT\_RETURN (Ergebnis empfangen)   & \texttt{servicer.py} \\
\#17 & GET\_RESULT (Ergebnis an Client)      & \texttt{servicer.py} \\
\#18 & Timeout + Retry-Logik                 & \texttt{dispatch\_loop.py} \\
\#19 & Strukturiertes Logging                & \texttt{src/common/logger.py} \\
\#20 & GET\_STATUS Monitoring                & \texttt{http\_status\_server.py} \\
\#21 & Queue + Nebenläufigkeit               & \texttt{task\_queue.py}, \texttt{task\_store.py} \\
\bottomrule
\end{tabular}
\caption{Übersicht implementierter Dispatcher-Issues}
\end{table}

\subsection{Containerisierung (Issues \#34--\#38)}

Alle vier Dockerfiles und die docker-compose.yml wurden vollständig
implementiert. Details siehe Abschnitt~\ref{sec:technologien}.

\begin{table}[H]
\centering
\begin{tabular}{lll}
\toprule
\textbf{Issue} & \textbf{Titel} & \textbf{Ergebnis} \\
\midrule
\#34 & Dockerfile für Client     & \texttt{client/Dockerfile} \\
\#35 & Dockerfile für Dispatcher & \texttt{Dockerfile} (Root) \\
\#36 & Dockerfile für Worker     & \texttt{worker/Dockerfile} + \texttt{worker/entrypoint.sh} \\
\#37 & Dockerfile für Namensdienst & \texttt{nameservice/Dockerfile} \\
\#38 & docker-compose.yml        & \texttt{docker-compose.yml} \\
\bottomrule
\end{tabular}
\caption{Übersicht implementierter Docker-Issues}
\end{table}
